% ═══════════════════════════════════════════════════════════════════════════
% MATERIALÜBERSICHT: Beobachtungszeitraum Januar/Februar 2026
% Spanisch Spätbeginner (Jahrgangsübergreifend Kl. 10-12)
% ═══════════════════════════════════════════════════════════════════════════

\documentclass[11pt, a4paper]{article}

\usepackage[utf8]{inputenc}
\usepackage[T1]{fontenc}
\usepackage[ngerman]{babel}
\usepackage{helvet}
\renewcommand{\familydefault}{\sfdefault}

\usepackage[margin=2cm]{geometry}
\usepackage[table]{xcolor}
\usepackage{array}
\usepackage{tabularx}
\usepackage{booktabs}
\usepackage{longtable}
\usepackage{enumitem}
\usepackage{tikz}
\usepackage{fancyhdr}
\usepackage{lastpage}
\usepackage[hidelinks]{hyperref}
\usepackage{amssymb}

% Farben
\definecolor{spanisch}{HTML}{c41e3a}
\definecolor{grau}{HTML}{666666}
\definecolor{hellgrau}{HTML}{f5f5f5}
\definecolor{success}{HTML}{2a7a4b}
\definecolor{warning}{HTML}{b35c00}
\definecolor{gruppeA}{HTML}{2a7a4b}
\definecolor{gruppeB}{HTML}{1565c0}

% Kopf- und Fußzeile
\pagestyle{fancy}
\fancyhf{}
\renewcommand{\headrulewidth}{0.4pt}
\lhead{\small Materialübersicht -- Spanisch Spätbeginner}
\rhead{\small Beobachtungszeitraum Jan/Feb 2026}
\cfoot{\small Seite \thepage\ von \pageref{LastPage}}

% Symbole
\newcommand{\ja}{\textcolor{success}{\textbf{$\checkmark$}}}
\newcommand{\nein}{\textcolor{grau}{--}}
\newcommand{\puffer}{\textcolor{warning}{\textbf{P}}}

\begin{document}

% ═══════════════════════════════════════════════════════════════════════════
% TITELSEITE
% ═══════════════════════════════════════════════════════════════════════════

\begin{center}
\vspace*{2cm}

{\Huge\bfseries\color{spanisch} Materialübersicht}

\vspace{0.5cm}

{\Large Beobachtungszeitraum Januar/Februar 2026}

\vspace{1cm}

{\large Spanisch Spätbeginner\\
Jahrgangsübergreifend Kl. 10, 11, 12\\
}

\vspace{2cm}

\begin{tikzpicture}
\draw[very thick, spanisch] (0,0) rectangle (12,3);
\node[font=\Large\bfseries] at (6,2.2) {Beobachtungstermin};
\node[font=\Huge\bfseries, spanisch] at (6,1) {Di 03.02.2026};
\node[font=\normalsize] at (6,0.3) {2. Hälfte (45 Min)};
\end{tikzpicture}

\vspace{2cm}

\begin{tabular}{ll}
\textbf{Kurs:} & 10 SuS (Kl. 10: 3, Kl. 11: 3, Kl. 12: 4) \\
\textbf{Wochenstunden:} & 4 (Mo + Di je 90 Min) \\
\textbf{Lehrwerk:} & A\_tope.com (Cornelsen) \\
\textbf{Niveaus:} & Gruppe A (A2) / Gruppe B (B1) \\
\end{tabular}

\vfill

{\small Stand: \today}

\end{center}

\newpage

% ═══════════════════════════════════════════════════════════════════════════
% ZEITPLAN
% ═══════════════════════════════════════════════════════════════════════════

\section*{Zeitplan (5 Wochen)}

\begin{center}
\begin{tikzpicture}[scale=0.9]
% Woche 1
\draw[thick] (0,0) rectangle (7,2);
\node[font=\bfseries] at (3.5,1.7) {WOCHE 1: 05.--06.01.};
\node[anchor=west, font=\small] at (0.2,1.1) {DS 1 (Mo): Bucharbeit parallel};
\node[anchor=west, font=\small] at (0.2,0.5) {DS 2 (Di): Fortführung + Übergang};

% Woche 2
\draw[thick] (8,0) rectangle (15,2);
\node[font=\bfseries] at (11.5,1.7) {WOCHE 2: 12.--13.01.};
\node[anchor=west, font=\small] at (8.2,1.1) {DS 3 (Mo): + Assessment-Ankündigung};
\node[anchor=west, font=\small] at (8.2,0.5) {DS 4 (Di): Letzte Bucharbeit};

% Woche 3
\draw[thick, spanisch] (0,-3) rectangle (7,-1);
\node[font=\bfseries, spanisch] at (3.5,-1.3) {WOCHE 3: 19.--20.01.};
\node[anchor=west, font=\small] at (0.2,-1.9) {DS 5 (Mo): \textbf{ASSESSMENT}};
\node[anchor=west, font=\small] at (0.2,-2.5) {DS 6 (Di): Brückenthema (OVERLAP)};

% Woche 4 (Puffer)
\draw[thick, warning, fill=yellow!10] (8,-3) rectangle (15,-1);
\node[font=\bfseries, warning] at (11.5,-1.3) {WOCHE 4: 26.--27.01. [PUFFER]};
\node[anchor=west, font=\small] at (8.2,-1.9) {DS 7 (Mo): Cultura y Diversión};
\node[anchor=west, font=\small] at (8.2,-2.5) {DS 8 (Di): Creatividad y Juego};

% Woche 5 (Beobachtung)
\draw[very thick, spanisch, fill=spanisch!10] (4,-6) rectangle (11,-4);
\node[font=\bfseries, spanisch] at (7.5,-4.3) {WOCHE 5: 02.--03.02.};
\node[anchor=west, font=\small] at (4.2,-4.9) {DS 9 (Mo): Vocabulario y Gramática};
\node[anchor=west, font=\small, spanisch] at (4.2,-5.5) {DS 10 (Di): \textbf{BEOBACHTUNG}};

% Pfeile
\draw[->, thick] (7,1) -- (8,1);
\draw[->, thick] (7.5,-0.5) -- (7.5,-1);
\draw[->, thick] (7,-2) -- (8,-2);
\draw[->, thick] (11.5,-3) -- (11.5,-4);
\draw[->, thick] (7.5,-3) -- (7.5,-4);
\end{tikzpicture}
\end{center}

\vspace{1cm}

\textbf{Legende:}
\begin{itemize}[leftmargin=2em]
    \item \textcolor{spanisch}{\textbf{Rot:}} Bewertungsrelevant (Assessment, Beobachtung)
    \item \textcolor{warning}{\textbf{Orange:}} Pufferwoche (entbehrlich, Online-Backup vorhanden)
    \item \textbf{Schwarz:} Regulärer Unterricht
\end{itemize}

\newpage

% ═══════════════════════════════════════════════════════════════════════════
% MATERIALMATRIX
% ═══════════════════════════════════════════════════════════════════════════

\section*{Materialmatrix}

\begin{center}
\small
\begin{tabular}{|c|l|c|c|c|c|c|}
\hline
\rowcolor{spanisch}
\textcolor{white}{\textbf{DS}} &
\textcolor{white}{\textbf{Datum / Thema}} &
\textcolor{white}{\textbf{LE}} &
\textcolor{white}{\textbf{AB}} &
\textcolor{white}{\textbf{Online}} &
\textcolor{white}{\textbf{Material}} &
\textcolor{white}{\textbf{Score}} \\
\hline
\hline
1 & Mo 05.01. -- Bucharbeit & \ja & \nein & \nein & \nein & 9.2 \\
\hline
2 & Di 06.01. -- Übergang & \ja & \nein & \nein & \nein & 9.1 \\
\hline
3 & Mo 12.01. -- Assessment-Ankündigung & \ja & \nein & \nein & \nein & 9.2 \\
\hline
4 & Di 13.01. -- Letzte Bucharbeit & \ja & \nein & \nein & \nein & 9.0 \\
\hline
\rowcolor{spanisch!10}
5 & Mo 19.01. -- \textbf{ASSESSMENT} & \ja & \nein & \nein & \nein & 9.3 \\
\hline
\rowcolor{spanisch!10}
6 & Di 20.01. -- Brückenthema (OVERLAP) & \ja & \nein & \nein & \nein & 9.6 \\
\hline
\rowcolor{warning!10}
7 & Mo 26.01. -- Cultura [PUFFER] & \ja & \nein & \ja & \nein & 9.1 \\
\hline
\rowcolor{warning!10}
8 & Di 27.01. -- Creatividad [PUFFER] & \ja & \nein & \ja & \nein & 9.2 \\
\hline
9 & Mo 02.02. -- Einstimmung & \ja & \nein & \ja & \nein & 9.5 \\
\hline
\rowcolor{spanisch!10}
10 & Di 03.02. -- \textbf{BEOBACHTUNG} & \ja & \nein & \nein & \ja & 9.5 \\
\hline
\end{tabular}
\end{center}

\vspace{0.5cm}

\textbf{Legende:}
\begin{itemize}[leftmargin=2em]
    \item \textbf{LE} = Langentwurf (Tier-1, vollständig)
    \item \textbf{AB} = Arbeitsblatt (Print)
    \item \textbf{Online} = Online-Aufgaben (bei Lehrerausfall)
    \item \textbf{Material} = Zusatzmaterial (Karten, Poster, etc.)
    \item \textbf{Score} = Didaktik-Score (min. 9.0 für Freigabe)
\end{itemize}

\newpage

% ═══════════════════════════════════════════════════════════════════════════
% DATEILISTE
% ═══════════════════════════════════════════════════════════════════════════

\section*{Dateiliste}

\subsection*{Langentwürfe (10 Dateien)}

\begin{tabular}{lll}
\toprule
\textbf{Datei} & \textbf{Größe} & \textbf{Inhalt} \\
\midrule
LE-01-bucharbeit.pdf & 289 KB & DS 1: Parallele Bucharbeit A/B \\
LE-02-uebergang.pdf & 330 KB & DS 2: Fortführung + Übergang \\
LE-03-assessment-ankuendigung.pdf & 287 KB & DS 3: Assessment-Ankündigung \\
LE-04-letzte-bucharbeit.pdf & 98 KB & DS 4: Letzte Bucharbeit + Kriterien \\
LE-05-assessment.pdf & 272 KB & DS 5: Assessment (sonstige Note) \\
LE-06-brueckenthema.pdf & 338 KB & DS 6: Brückenthema (1. Overlap) \\
LE-07-puffer-cultura.pdf & 329 KB & DS 7: Cultura [Puffer] \\
LE-08-puffer-creatividad.pdf & 330 KB & DS 8: Creatividad [Puffer] \\
LE-09-einstimmung.pdf & 340 KB & DS 9: Vocabulario y Gramática \\
LE-10-beobachtung.pdf & 119 KB & DS 10: Beobachtungsstunde \\
\bottomrule
\end{tabular}

\vspace{1cm}

\subsection*{Zusatzmaterial (1 Datei)}

\begin{tabular}{lll}
\toprule
\textbf{Datei} & \textbf{Größe} & \textbf{Inhalt} \\
\midrule
MATERIAL-10-fragekarten.pdf & 54 KB & Fragekarten A (6x), B (6x), Zusammenfassung (4x), \\
 & & Tandem-Übersicht, Vokabel-Poster \\
\bottomrule
\end{tabular}

\newpage

% ═══════════════════════════════════════════════════════════════════════════
% STUNDEN-DETAILS
% ═══════════════════════════════════════════════════════════════════════════

\section*{Stundenübersicht im Detail}

\subsection*{DS 1 -- Mo 05.01.2026: Bucharbeit parallel}

\begin{tabular}{ll}
\textbf{Thema:} & Parallele Bucharbeit nach Niveaus \\
\textbf{Gruppe A:} & A\_tope.com S. 30--31 (Mi barrio) \\
\textbf{Gruppe B:} & A\_tope.com S. 102--103 (El mundo laboral) \\
\textbf{Ziel:} & Beide Gruppen arbeiten selbstständig an ihrem Niveau \\
\textbf{Materialien:} & Langentwurf, Lehrwerk \\
\end{tabular}

\vspace{0.5cm}

\subsection*{DS 2 -- Di 06.01.2026: Fortführung + Übergang}

\begin{tabular}{ll}
\textbf{Thema:} & Fortsetzung der Bucharbeit \\
\textbf{Gruppe A:} & A\_tope.com S. 32--33 \\
\textbf{Gruppe B:} & A\_tope.com S. 104--105 \\
\textbf{Ziel:} & Abschluss der aktuellen Lektionen \\
\textbf{Materialien:} & Langentwurf, Lehrwerk \\
\end{tabular}

\vspace{0.5cm}

\subsection*{DS 3 -- Mo 12.01.2026: Assessment-Ankündigung}

\begin{tabular}{ll}
\textbf{Thema:} & Bucharbeit + Ankündigung Assessment \\
\textbf{Gruppe A:} & A\_tope.com S. 34--35 \\
\textbf{Gruppe B:} & A\_tope.com S. 106--107 \\
\textbf{Ziel:} & Transparenz über anstehendes Assessment \\
\textbf{Materialien:} & Langentwurf, Lehrwerk \\
\end{tabular}

\vspace{0.5cm}

\subsection*{DS 4 -- Di 13.01.2026: Letzte Bucharbeit}

\begin{tabular}{ll}
\textbf{Thema:} & Abschluss Bucharbeit + Kriterienvergabe \\
\textbf{Gruppe A:} & Punto final (E-Mail schreiben) \\
\textbf{Gruppe B:} & Evaluación 1 \\
\textbf{Ziel:} & Vorbereitung auf Assessment, Kriterien bekannt \\
\textbf{Materialien:} & Langentwurf, Bewertungsbögen (im LE) \\
\end{tabular}

\vspace{0.5cm}

\subsection*{DS 5 -- Mo 19.01.2026: ASSESSMENT}

\begin{tabular}{ll}
\textbf{Thema:} & \textbf{Sonstige Leistung (mündlich)} \\
\textbf{Format:} & Rotation: Gruppe A präsentiert, B arbeitet still (dann Wechsel) \\
\textbf{Bewertung:} & 18-Punkte-Schema → 14 NP \\
\textbf{Ziel:} & Bewertung vor der Beobachtung abgeschlossen \\
\textbf{Materialien:} & Langentwurf, Bewertungsbögen (Anlage im LE) \\
\end{tabular}

\vspace{0.5cm}

\subsection*{DS 6 -- Di 20.01.2026: Brückenthema (ERSTER OVERLAP)}

\begin{tabular}{ll}
\textbf{Thema:} & \textbf{Personas y características} \\
\textbf{Besonderheit:} & Erster gemeinsamer Inhalt für A und B \\
\textbf{Format:} & Experten-Interview (B interviewt A) \\
\textbf{Ziel:} & Vorbereitung auf gemischte Tandems in Beobachtung \\
\textbf{Materialien:} & Langentwurf \\
\end{tabular}

\newpage

\subsection*{DS 7 -- Mo 26.01.2026: Cultura y Diversión [PUFFER]}

\begin{tabular}{ll}
\textbf{Thema:} & Kulturelle Inhalte + Spielerische Festigung \\
\textbf{Charakter:} & \textcolor{warning}{\textbf{PUFFER}} -- entbehrlich für Beobachtung \\
\textbf{Präsenz:} & Música, Quiz-Rallye, Adjektiv-Bingo, Personenraten \\
\textbf{Online:} & LearningApps, Quizlet, Video + Fragen \\
\textbf{Materialien:} & Langentwurf, AB-07-A/B \\
\end{tabular}

\vspace{0.5cm}

\subsection*{DS 8 -- Di 27.01.2026: Creatividad y Juego [PUFFER]}

\begin{tabular}{ll}
\textbf{Thema:} & Kreative Produktion + Spiele \\
\textbf{Charakter:} & \textcolor{warning}{\textbf{PUFFER}} -- entbehrlich für Beobachtung \\
\textbf{Präsenz:} & Wortschatz-Staffel, Mini-Poster, Gallery Walk, Escape-Room \\
\textbf{Online:} & Schreibaufgabe, Selbsttest, Padlet \\
\textbf{Materialien:} & Langentwurf, AB-08-A/B \\
\end{tabular}

\vspace{0.5cm}

\subsection*{DS 9 -- Mo 02.02.2026: Vocabulario y Gramática}

\begin{tabular}{ll}
\textbf{Thema:} & Wiederholung ohne Neues \\
\textbf{Charakter:} & Bekannte Inhalte in neuer Kombination (Gamification) \\
\textbf{Präsenz:} & Repaso Relámpago, Lernpfad, Speed-Dating, Tandem-Challenge \\
\textbf{Online:} & Lernpfad + Schreibaufgabe + Selbst-Checkliste \\
\textbf{Materialien:} & Langentwurf, AB-09-A/B \\
\end{tabular}

\vspace{0.5cm}

\subsection*{DS 10 -- Di 03.02.2026: BEOBACHTUNG}

\begin{tabular}{ll}
\textbf{Thema:} & \textbf{Personas y características -- Interview-Karussell} \\
\textbf{Beobachtet:} & 2. Hälfte (45 Min) \\
\textbf{Phase 1:} & Homogene Paare (A$\leftrightarrow$A, B$\leftrightarrow$B) -- 12 Min \\
\textbf{Phase 2:} & Gemischte Tandems, Ankreuz-Vorlage -- 18 Min \\
\textbf{Phase 3:} & Präsentation (nur 3 Tandems) -- 9 Min \\
\textbf{Materialien:} & Langentwurf, MATERIAL-10-fragekarten.pdf \\
\end{tabular}

\newpage

% ═══════════════════════════════════════════════════════════════════════════
% CHECKLISTE VOR DER BEOBACHTUNG
% ═══════════════════════════════════════════════════════════════════════════

\section*{Checkliste vor der Beobachtung}

\subsection*{Materialien drucken/vorbereiten}

\begin{itemize}[leftmargin=2em]
    \item[$\square$] Fragekarten A (6x) ausdrucken und laminieren
    \item[$\square$] Fragekarten B (6x) ausdrucken und laminieren
    \item[$\square$] Zusammenfassungs-Vorlage (5x) ausdrucken
    \item[$\square$] Vokabel-Poster (1x) groß ausdrucken für Tafel
    \item[$\square$] Tandem-Übersicht (1x) für eigene Referenz
    \item[$\square$] Timer bereitstellen (sichtbar für alle)
\end{itemize}

\subsection*{Inhaltliche Vorbereitung}

\begin{itemize}[leftmargin=2em]
    \item[$\square$] Tandem-Paarungen bekannt (H.P.$\leftrightarrow$E.H., A.A.$\leftrightarrow$V.S., etc.)
    \item[$\square$] Fokus-SuS identifiziert (H.P., A.A., A.S.)
    \item[$\square$] Redemittel für Zusammenfassung geübt (Se llama..., Es..., Le gusta...)
    \item[$\square$] ser vs. estar gefestigt
\end{itemize}

\subsection*{Am Tag der Beobachtung}

\begin{itemize}[leftmargin=2em]
    \item[$\square$] Raum vorbereitet (Tische für Partnerarbeit)
    \item[$\square$] Fragekarten sortiert und bereitgelegt
    \item[$\square$] Timer getestet
    \item[$\square$] Vokabel-Poster aufgehängt
    \item[$\square$] Langentwurf LE-10 für Beobachter kopiert
\end{itemize}

\vspace{1cm}

\begin{center}
\fbox{%
\begin{minipage}{0.8\textwidth}
\centering
\vspace{0.3cm}
{\Large\bfseries ¡Mucha suerte!}\\[0.3cm]
{\normalsize Du bist bestens vorbereitet.}\\[0.2cm]
{\small Die SuS kennen alle Formate. Du kennst deine SuS.\\
Vertraue dem Prozess.}
\vspace{0.3cm}
\end{minipage}
}
\end{center}

\end{document}
